\chapter{Related Work}

% what do existing reviews cover and what is still missing
Existing reviews and conceptual frameworks have made important contributions to organizing the growing body of research on NAS in FL. Their classifications explain broad procedural differences between approaches and, more recently, the design decisions from which these differences arise. However, their primary unit of analysis is the complete NAS-in-FL method or its design configuration. The challenges that motivate individual design decisions and the solution mechanisms shared across methods consequently remain distributed throughout the primary literature.

% explain what earlier reviews have done
Earlier reviews organize NAS-in-FL research around broad procedural and optimization categories. One review distinguishes online from offline implementations and single-objective from multi-objective searches, thereby showing when architecture search occurs and which objectives guide it~\cite{fl_to_nas_survey_2021}. A broader review compares NAS and hyperparameter optimization in centralized and federated environments and summarizes how individual approaches modify conventional NAS components for FL~\cite{nas_hpo_fl_survey_2023}. Research on multi-objective methods in FL offers an adjacent perspective by considering architecture search where it contributes to the joint optimization of predictive performance and system objectives~\cite{multi_objective_methods_in_fl_2025}. These classifications provide an overview of the available approaches, but they do not consolidate the recurring challenges and solutions embedded within them.

% relation to Nick's paper
More recent work provides a finer-grained account of NAS-in-FL design. A unified design concept decomposes the NAS-in-FL pipeline into design dimensions, identifies recurring method archetypes and relates their configurations to method traits and suitable FL systems~\cite{nas_in_fl_2026}. This representation supplies a common vocabulary for comparing approaches that use different terminology and distribute search activities differently between clients and the server. Its analytical focus remains the configuration of complete approaches, leaving the relationships between particular challenges and reusable solution mechanisms implicit within those configurations.

% explain gap
Reviews of NAS and FL separately provide concepts needed to understand these relationships. NAS research distinguishes the search space, search strategy and performance-estimation strategy~\cite{nas_survey_2019}, while FL taxonomies organize concerns such as heterogeneity, efficiency, security and aggregation~\cite{fl_taxonomy_2024}. Each describes one side of the combined setting. In NAS-in-FL, their interaction can change the problem. For example, the cost of candidate evaluation becomes a distributed client workload, while differences between searched architectures can invalidate the shared representation assumed by conventional aggregation~\cite{nas_in_fl_2026,rafl_2024}. These adjacent reviews therefore establish relevant context without providing an account of the challenge and solution space created by their combination.

% explain contribution
This thesis complements the existing classifications by using recurring challenges and solutions as its units of synthesis. It consolidates descriptions that occur under different terminology, distinguishes solutions by their mechanisms and maps the relationships between them. The analysis also retains the FL system characteristics and NAS designs that delimit where a challenge arises and where a solution is applicable. This concept-centric perspective explains why particular changes to conventional NAS are needed and makes visible how the same solution can recur across otherwise different NAS-in-FL methods.
